\begingroup
\renewcommand{\arraystretch}{1.45}
\InterpolationArticle{17}
It follows readily from algebraic principles that, if, in the first case of the preceding article, the product
\[
(\cos t-\cos a)(\cos t-\cos b)(\cos t-\cos c)(\cos t-\cos d)\text{ etc.}
\]
and, in the second case, the product
\[
\sin t(\cos t-\cos a)(\cos t-\cos b)(\cos t-\cos c)(\cos t-\cos d)\text{ etc.}
\]
is denoted by \(P\), every function similar to \(X\) that satisfies the same values
\(A,B,C,D\) etc. for \(x=a,b,c,d\) etc. must necessarily be contained in the form \(X+PY\),
where \(Y\) is an undetermined function of the arc \(x\) of the form
\[
\gamma+\gamma'\cos x+\gamma''\cos 2x+\gamma'''\cos 3x+\text{etc.}
\]
The converse theorem, namely that every function contained in the form \(X+PY\) satisfies those values, is self-evident. Moreover, it is evident that the product
\(P\) is of the \(m+1^{\mathrm{th}}\) order, and thus of the order immediately above that of \(X\); hence no other function
similar to \(X\) that satisfies the same values exists except one of higher order
than \(X\).

Conversely, if a function \(X'\), similar to \(X\) but of higher order,
is known from elsewhere and satisfies the given values \(A,B,C,D\) etc., the function \(Y\),
of the form we have given, can be determined so that \(X'+PY\) is reduced to the
\endgroup
\begingroup
\renewcommand{\arraystretch}{1.45}
\SourceLeaf{302}
\noindent
\(m^{\mathrm{th}}\) order; this must necessarily yield the very function \(X\) obtained from the values \(A,B,C,D\) etc.
by the method of the preceding article. Since the method of finding the highest term
of the function \(Y\) and then, successively, the lower terms can be derived as a special case
of the method given in Article 15, II, there is no need to dwell on this matter
any longer here.

\InterpolationArticle{18}
All the preceding investigations rest on the condition that \(X\) is a series ending with the cosine and sine of the arc \(mx\). If this condition does not
hold, as many values of \(X\) as we have so far assumed known clearly do not
suffice for a complete determination, and the expression for \(T\) given in the various cases we have considered will need a correction that depends on the following coefficients
in \(X\) and can readily be specified by the summations of Articles 1 and 2. If, however,
the series \(X\), whether it continues to infinity or ends somewhere, converges to such a degree
that the parts following \(\cos mx\) and \(\sin mx\) can be treated as vanishing,
that correction may also be neglected, and we will thus have obtained a value of \(T\) that is at least approximately
true. However this may be, \(T\) certainly represents the simplest function similar
to the true one by which all the given values can be satisfied.

\endgroup
